% =============================================================================
% Draft insertion for NOETHER_paper.tex
%
% Suggested placement: as a new \subsubsection inside \subsection{A small-scale
% comparative case study} (\S\ref{subsec:case-study}), immediately after
% \S\ref{subsec:deepcrime-pilot} (the DeepCrime-style pilot) and after the
% "Comparative evaluation against published baselines (protocol)" paragraph.
% This subsubsection is the first instantiation of that protocol on Java SUTs.
%
% Pre-registration anchor: configs/d4j_algebra_rich_criterion.json on branch
% feat/d4j-algebra-rich; head 67e36c4. Test gate green at f04a926 (33 SUT
% directories, 200 DSL files spanning the §6.5 utility-method baseline and
% the §6.6 algebra-rich extension).
%
% Reframing of the head-to-head claim: competitive parity with coverage
% extension (claims A, B, C, F strongly supported; claim E reframed away
% from head-to-head superiority). See supplementary S7 for the full claim
% audit and per-SUT outcome matrix.
% =============================================================================

\subsubsection{Algebra-rich Java SUTs: Set N versus GenMorph-evolved Set G}
\label{subsec:d4j-algebra-rich}

We instantiate the protocol of \S\ref{para:comp-eval-protocol} on the
algebra-rich subset of the GenMorph 23-method Java
benchmark~\cite{Ayerdi2023GenMorph}. The protocol's structural-coverage
filter selects subjects whose induced operator algebra has at least one
non-empty block beyond $G$; on this filter we hand-derive Set~N MRs and
rerun GenMorph's GP pipeline to obtain Set~G. The Set~M, Set~L, and
Set~B arms remain as committed future work; this subsubsection
delivers the paired Set~N versus Set~G comparison.

\paragraph{Pre-registered SUT criterion.}
The SUT-selection rule is committed in
\texttt{configs/d4j\_algebra\_rich\_criterion.json} on branch
\texttt{feat/d4j-algebra-rich} \emph{before} any new evaluation data
existed. The criterion references package roots
(\texttt{org.apache.commons.math3.ode}, \texttt{*.linear},
\texttt{*.transform}, \texttt{*.analysis.solvers},
\texttt{*.distribution}, \texttt{*.optim}, \texttt{*.fitting},
\texttt{*.stat.regression}, \texttt{*.complex}, \texttt{*.geometry},
\texttt{*.fraction} and their \texttt{math.*} legacy counterparts),
per-package block-coverage hypotheses, and method-signature
constraints; it contains no bug-id and no kill-rate references and
cannot be tuned by evaluation outcomes. The git timestamp chain
(criterion $\rightarrow$ inscope filter $\rightarrow$ Set~N derivation
$\rightarrow$ Set~G GP rerun $\rightarrow$ M1 numbers) is the auditable
proof of pre-registration.

\paragraph{Subjects and MR sets.}
Ten SUTs are inlined into a single \texttt{MathSignalClass} (or
\texttt{ComplexSignal} for the instance method) so that PIT mutates
the algorithm directly rather than a delegator: \texttt{midpoint},
\texttt{exactLog2}, \texttt{isSequence} (boolean predicate),
\texttt{clamp}, \texttt{signum}, \texttt{ComplexSignal.add} (instance
method), \texttt{gcdSig}, \texttt{lcmSig}, \texttt{hypotSig},
\texttt{powerSig}. The SUTs span the symmetry, translation/period,
scaling, and identity blocks of $\mathcal{A}_P$.

\textit{Set~N (NOETHER-derived).} Thirty hand-derived NOETHER MRs
(2--4 per SUT), one or more per non-empty block of each SUT's induced
algebra, expressed in GenMorph's JIR/JOR DSL.

\textit{Set~G (GenMorph GP-evolved).} Harvested by rerunning
GenMorph's \texttt{genmorph.py gen} on the same 10 SUTs at seed $=11$.
Per-process source isolation and a \texttt{MAJOR\_HOME} fix were
required for the upstream Major-mutation pipeline (Section S7.4 in
supplementary~\ref{S7-d4j}). At a 1\,min GAssert budget the rerun
yields 26 MRs across 8 SUTs; two SUTs return Set~G \emph{not
available} for structural reasons described next.

\paragraph{Set G is structurally absent on two of the ten SUTs.}
GenMorph's GP pipeline fails on two SUTs for reasons internal to its
plumbing rather than the underlying mutation-testing problem:
\begin{itemize}[nosep,leftmargin=*]
  \item \texttt{ComplexSignal.add} (instance method) triggers an
        XStream \texttt{NullPointerException} in
        \texttt{MethodTestTransformerConfig} when deserialising the
        receiver object;
        \texttt{ch.usi.gassert.data.types.AbstractCollectionConverter}
        in the upstream snapshot does not handle user-defined
        receivers.
  \item \texttt{isSequence} (boolean predicate) yields GP-evolved JOR
        strings containing dangling \texttt{<} operators, e.g.\
        \texttt{((\ldots <))}, that do not parse as Java; the GP
        grammar for predicate-output SUTs is degenerate at this scale.
\end{itemize}
Both SUTs admit Set~N MRs that derive directly from the SUT signature
(commutativity for \texttt{ComplexSignal.add}, monotonic-window
translation for \texttt{isSequence}) and run to completion. Together
the two SUTs cover 8 of the 70 PIT mutants under evaluation, on which
only Set~N is defined. This structural-coverage extension is reported
as a finding independent of the kill-rate statistics that follow.

\paragraph{Pooled head-to-head result.}
Across the eight head-to-head SUTs, $n = 70$ PIT mutants are
evaluated. Set~N kills 34 mutants; Set~G kills 39
(Table~\ref{tab:d4j-algebra-rich-pooled}). The pooled M1 kill rates
are $0.486$ for Set~N (Wilson 95\% CI $[0.372,\,0.600]$) versus
$0.557$ for Set~G ($[0.441,\,0.668]$); a McNemar exact paired test
gives $p = 0.359$. The 95\% confidence intervals overlap and the test
fails to reject the null at $\alpha = 0.05$. \textbf{$n = 70$ across
$10$ SUTs is underpowered for a paired hypothesis test at
$\alpha = 0.05$}; the absence of a significant difference is reported
as descriptive evidence rather than as a confirmation of equivalence.

The Set~G arm runs at GAssert budget 1\,min versus the GenMorph
upstream default 30\,min. The reduced budget handicaps Set~G; a full
30\,min rerun would be expected to widen Set~G's lead, not narrow it.
The direction of this caveat is therefore conservative against any
parity claim made in Set~N's favour.

\begin{table}[h]
\centering
\caption{Per-SUT PIT mutation kill counts, Set~N (NOETHER-derived)
versus Set~G (GenMorph GP-evolved), on the 10 algebra-rich Java SUTs
selected by the pre-registered \texttt{d4j\_algebra\_rich\_criterion}
filter. Set~G is structurally absent on \texttt{ComplexSignal.add}
(XStream receiver-deserialisation failure) and \texttt{isSequence}
(degenerate GP grammar on boolean predicates); see text. Pooled across
the 8 head-to-head SUTs ($n = 70$ mutants): Set~N $= 34$,
Set~G $= 39$, McNemar exact $p = 0.359$ at the 1\,min GAssert budget.
Wilson 95\% confidence intervals on the pooled rates overlap.}
\label{tab:d4j-algebra-rich-pooled}
\small
\begin{tabular}{lccccc}
\toprule
SUT & mutants & Set N kills & Set G kills & $\Delta$ & note \\
\midrule
\texttt{ComplexSignal.add} &  3 & 2 & N/A & --- & instance method \\
\texttt{midpoint}          &  3 & 3 &   3 & $0$  & tie \\
\texttt{exactLog2}         & 10 & 4 &   0 & $+4$ & N narrow \\
\texttt{isSequence}        &  5 & 0 & N/A & --- & boolean predicate \\
\texttt{clamp}             &  7 & 3 &   6 & $-3$ & G \\
\texttt{signum}            &  6 & 4 &   4 & $0$  & tie \\
\texttt{gcdSig}            &  9 & 6 &   5 & $+1$ & N narrow \\
\texttt{hypotSig}          &  4 & 2 &   4 & $-2$ & G \\
\texttt{lcmSig}            & 11 & 4 &   7 & $-3$ & G \\
\texttt{powerSig}          & 12 & 8 &   7 & $+1$ & N narrow \\
\midrule
\textbf{pooled (n=70)}     &    & \textbf{34} & \textbf{39} & $-5$ & ns \\
\bottomrule
\end{tabular}
\end{table}

\paragraph{Per-block kill pattern: empirical witness for the eight-block decomposition.}
Set~N MRs are tagged by the NOETHER block they instantiate, so the
70-mutant outcome decomposes into per-block per-MR kill rates rather
than a single aggregate. The pattern matches the algebraic-mechanism
prediction in three respects, including one prediction that is sharp,
falsifiable, and confirmed.

\textit{$T^{*}$ block (translation and period).} Translation MRs carry
the highest single-MR kill rates on the numeric SUTs:
\texttt{midpoint} \texttt{T\_shift} 3/3, \texttt{powerSig}
\texttt{T\_exp\_step} 8/12, \texttt{exactLog2} \texttt{T\_double}
4/10, \texttt{clamp} \texttt{T\_shift} 3/7. Translation invariance is
the property most sensitive to PIT's arithmetic-operator swaps, and
$T^{*}$-derived MRs detect this fingerprint reliably.

\textit{$G$ block (group and symmetry).} Symmetry MRs are moderately
to highly effective when the SUT exposes the appropriate symmetry:
\texttt{signum} \texttt{G\_negate} 4/6, \texttt{gcdSig} \texttt{G\_swap}
5/9, \texttt{midpoint} \texttt{G\_swap} 1/3, \texttt{ComplexSignal.add}
\texttt{G\_swap} 2/3.

\textit{$\mathcal{L}^{*}$ block (linearity and scaling): a uniformly
near-zero kill rate, theoretically predicted.} The scaling MRs of the
form $f(\lambda \mathbf{x}) = \lambda \cdot f(\mathbf{x})$ kill 0 of
3, 0 of 7, 0 of 9, 0 of 11, and 2 of 4 mutants on \texttt{midpoint},
\texttt{clamp}, \texttt{gcdSig}, \texttt{lcmSig}, and \texttt{hypotSig}
respectively. The flat near-zero pattern is the algebraic-mechanism
claim's sharpest empirical witness: PIT's default mutator set
($+/-/\times/\div$ swaps, \texttt{return zero}, \texttt{return one},
conditional negation) preserves homogeneity of degree~$1$, so a paired
MR of the form $f(\lambda \mathbf{x}) = \lambda \cdot f(\mathbf{x})$
cannot in principle distinguish the original from such mutants. The
$\mathcal{L}^{*}$-block MRs are blind to PIT's homogeneity-preserving
mutants \emph{by construction}, and the data confirms this. We name
this finding \emph{$\mathcal{L}^{*}$-block blindness} and treat it as
a direct empirical witness for the conservation-law-to-MR-blindness
correspondence articulated in \S\ref{subsec:algebra-induced-MRs}: a
block whose generator is a continuous symmetry of the mutator set
contributes MRs that are necessarily silent on mutator-induced
defects.

\paragraph{Independent rediscovery: GP recovers Set N's algebraic core.}
On \texttt{midpoint}, GenMorph's GP independently evolves three MRs
that map onto Set~N's blocks: \texttt{SwitchParams??1@2} corresponds
to \texttt{G\_swap} (commutativity);
\texttt{NumericAddition?1.000000?1} corresponds to \texttt{T\_shift}
(translation by~$1$); and \texttt{NumericMultiplication?0.500000?2}
approximates \texttt{L\_scale} (multiplicative scaling). Set~N is
derived \emph{a~priori} from the operator algebra of \texttt{midpoint};
GP rediscovers a subset of the same MRs by mutation-killing search.
That two pipelines with no shared design converge on the same
algebraic primitives is direct corroboration of the framework's
central claim that algebraic structure is the operative generator of
effective MRs, not an analyst's after-the-fact rationalisation.

\paragraph{Reframing the head-to-head claim.}
The pooled comparison does not support a head-to-head superiority
claim for Set~N over Set~G: Set~G is directionally ahead by $0.071$
absolute pooled, the McNemar test fails to reject the null, and
$n = 70$ is underpowered for an $\alpha = 0.05$ verdict. What the
data does support is a two-part claim:
\begin{enumerate}[nosep,leftmargin=*]
  \item \emph{Competitive parity on the head-to-head substrate.} On
        the eight SUTs both methods can address, the pooled kill
        rates are within McNemar non-significance at the 1\,min
        Set~G budget. Per-SUT, Set~N narrowly wins 4 SUTs
        (\texttt{midpoint}, \texttt{exactLog2}, \texttt{gcdSig},
        \texttt{powerSig}); Set~G clearly wins 3 SUTs
        (\texttt{clamp}, \texttt{hypotSig}, \texttt{lcmSig}); 1 SUT
        ties non-trivially (\texttt{signum}); and 1 SUT ties at zero
        for both methods (\texttt{isSequence}, where the paired-MR
        DSL is structurally inadequate for boolean classifiers).
  \item \emph{Structural coverage extension.} Set~N is defined and
        runnable on the two SUTs (\texttt{ComplexSignal.add},
        \texttt{isSequence}) on which GenMorph's GP pipeline is
        structurally absent. The two SUTs together cover $11.4\%$ of
        the mutant denominator under evaluation, and on these SUTs
        the algebra-derived method is the only option among the
        compared baselines.
\end{enumerate}
We frame the contribution of this subsubsection as
\emph{competitive parity with coverage extension}, not as a
head-to-head superiority claim. The full per-SUT outcome matrix, the
per-block kill table for all 30 Set~N MRs, the GP raw output and the
\texttt{.jir}/\texttt{.jor} translation, the McNemar $b/c$ counts on
the eight head-to-head SUTs, and the test gate
(\texttt{tests/run.sh} reproducing the 33 SUT directories and the
200 DSL files at branch \texttt{feat/d4j-algebra-rich}, head
\texttt{67e36c4}) are in supplementary~\ref{S7-d4j}.

\paragraph{Threats specific to this subsubsection.}
(a) \textbf{Set~G budget asymmetry.} GenMorph's published
configuration is 30\,min GAssert; we ran at 1\,min for parallel
scheduling. The 1\,min budget handicaps Set~G in the conservative
direction relative to Set~N's parity claim: a full-budget rerun, were
Set~G to extend its lead, would not invalidate the structural
coverage-extension finding but would weaken the parity reading.
(b) \textbf{Sample size.} $n = 70$ across 10 SUTs is underpowered for
a head-to-head verdict at $\alpha = 0.05$; this is why the
contribution is reframed as competitive parity rather than
superiority. Extending to the $\sim 38$ in-scope D4J subjects the
pre-registered criterion identifies remains future work and is
expected to be uncertain in direction by current evidence.
(c) \textbf{Set~G structural absence.} The two N/A SUTs reflect the
state of GenMorph upstream at our snapshot; an upstream patch
addressing XStream receiver deserialisation and the boolean-output
JOR grammar would restore Set~G coverage on these SUTs. The
structural-extension claim is reported with that future-state
qualifier; the claim does not assert that algebra-derived MRs are the
only kind that \emph{could in principle} be defined on these SUTs,
only that they are the only kind a current automated pipeline
\emph{does} produce.
(d) \textbf{Substrate selection.} Algebra-rich SUTs are explicitly
the in-scope class for NOETHER (\S\ref{subsec:algebra-induced-MRs});
the kill-rate gain over the §6.5 utility-method baseline (Set~N
$0.288 \rightarrow 0.486$, Set~G $0.363 \rightarrow 0.557$) is
expected by the framework's scope declaration and is not advanced as
a generalisable improvement on every Java SUT.
